\documentclass[11pt]{article}
\usepackage{amsmath, amssymb, geometry, enumitem}
\geometry{margin=1in}
\title{Legacy ScoreSaber Score Reconstruction}
\author{}
\date{}

\begin{document}
   \maketitle

   \section{What We Have}

   Legacy ScoreSaber replays provide per-frame tuples $(t_{i},\; s_{i},\; k_{i})$,
   namely timestamp, cumulative score, and combo. No individual note scores. We
   have to figure those out ourselves.

   \section{Scoring Refresher}

   Note score,
   \[
      B = B_{\text{pre}}+ B_{\text{post}}+ B_{\text{acc}}
   \]
   Pre swing up to 70, post swing up to 30, acc is 15. Max 115. Cumulative score
   multiplies each note by combo multiplier $m \in \{1, 2, 4, 8\}$.

   The catch is post swing doesn't get applied instantly. It settles over roughly
   50 to 100ms after the hit. So at the frame you hit a note the score only has
   \[
      s_{i} \approx S_{\text{prev}}+ (B_{\text{pre}}+ B_{\text{acc}}) \cdot m + \beta
      \cdot B_{\text{post}}\cdot m
   \]
   where $0 \le \beta \le 1$ is how much post swing has settled so far. The rest trickles
   in over the next few frames.

   \section{Finding Notes}

   We scan frames for combo changes. Combo goes up, hit. Combo drops, failed
   action. If combo jumps by more than 1 that means multiple notes were hit in the
   same frame.

   \section{Getting Scores Out}

   For a lone note with plenty of time before the next hit, just grab the score delta
   across the window and divide by the multiplier
   \[
      B = \text{clamp}\!\left(\left\lfloor \frac{s_{\text{end}}- s_{\text{before}}}{m}
      \right\rfloor,\; 0,\; 115\right)
   \]
   This works well when the window is wide enough for post swing to settle.

   \subsection{The Problem}

   When notes are close together (stacks, slider dots, doubles) the windows get
   tiny. The first note's window misses its post swing, the second note's window
   picks up the leftover. Second note gets clamped at 115, the excess is just
   gone. This adds up over the whole map.

   \section{Chaining}

   The fix is to chain consecutive hits that are closer than
   $\tau = 40\,\text{ms}$ apart. Use one big settled window for the whole chain,
   then split the total evenly.

   \subsection{Why 40ms}

   Gap between 1/4 stream notes,
   \[
      \Delta t = \frac{60}{4 \cdot \text{BPM}}\cdot 1000\;\text{ms}
   \]

   \begin{center}
      \renewcommand{\arraystretch}{1.2}
      \begin{tabular}{r r l}
         \hline
         BPM & Gap  & Chained?   \\
         \hline
         250 & 60ms & No         \\
         300 & 50ms & No         \\
         350 & 43ms & No         \\
         375 & 40ms & Borderline \\
         400 & 38ms & Yes        \\
         \hline
      \end{tabular}
   \end{center}

   Streams up to about 370 BPM won't get chained. Stacks (0ms) and slider dots (around
   18ms) will, which is the intended behavior.

   \subsection{Settled Window}

   For a chain from change $c$ to $c_{\text{end}}$,
   \[
      \Delta S = s_{f_{c_{\text{end}}} - 1}- s_{f_c - 1}
   \]
   Since $c_{\text{end}}$ is the first change after the chain with a big enough
   gap, by that point all the post swing from every note in the chain has settled.

   \subsection{Splitting the Score}

   Average base score using the multiplier sum,
   \[
      \bar{B}= \text{clamp}\!\left(\operatorname{round}\!\left(\frac{\Delta S}{\sum
      m_{j}}\right),\; 0,\; 115\right)
   \]

   \subsection{Adding Some Variation}

   Giving every note in a chain the exact same score looks artificial. So we
   nudge them using the immediate frame deltas as a guide. Let
   \[
      \delta_{g} = s_{f_g}- s_{f_g - 1}\qquad H = \sum_{g} h_{g}
   \]
   Then
   \[
      d_{g} = \frac{\max(1,\; \delta_{g})}{h_{g}}\qquad \bar{d}= \frac{1}{H}\sum d
      _{g} \cdot h_{g}
   \]
   \[
      \epsilon_{g} = \text{clamp}\!\left(\operatorname{round}\!\left(\left(\frac{d_{g}}{\max(1,\;
      \bar{d})}- 1\right) \cdot 15\right),\; -5,\; 5\right)
   \]

   For multi-hit changes where we only have one delta, just alternate $\{-1, +1\}$
   between notes. Final score,
   \[
      B = \text{clamp}\!\left(\bar{B}+ \epsilon_{g} + \omega,\; 0,\; 115\right)
   \]

   \section{Turning It Back into Cut Ratings}

   ScoreManager needs \texttt{beforeCutRating} and \texttt{afterCutRating}. We
   fill pre swing first, since that is what the game applies immediately.
   \[
      B_{\text{swing}}= \max(0,\; B - 15)
   \]
   \[
      r_{\text{pre}}= \min\!\left(1,\; \frac{B_{\text{swing}}}{70}\right) \qquad r
      _{\text{post}}=
      \begin{cases}
         \frac{B_{\text{swing}} - 70}{30} & \text{if }B_{\text{swing}}> 70 \\
         0                                & \text{otherwise}
      \end{cases}
   \]

   \section{How Well Does It Work}

   About 0.02\% off the real cumulative score. The remaining error is from
   clamping at 115 eating the excess when notes overlap, and from the fact that we
   just don't have per-note data in the legacy format. That limitation is
   inherent to the legacy format.
\end{document}